\documentclass[11pt]{article}
\input{style-sujet}

\newcommand{\note}[1]{\par\smallskip{\itshape\small #1}\par\smallskip}

\begin{document}

\entetesujet{Sujet bac 2017 -- Série C}

% ====================== EXERCICE 1 ======================
\exo{1}{4 points}

On considère l'équation $(E)$ dans $\Z \times \Z$ : $48x + 35y = 1$.

\begin{enumerate}[label=\alph*]
  \item Justifier à l'aide du théorème de Bézout que l'équation $(E)$ admet des solutions dans $\Z \times \Z$.
  \item Justifier que les entiers $48$ et $-8$ sont inverses modulo $35$.
  \item En remarquant que $(E)$ peut s'écrire $48x \equiv 1\ [35]$, déterminer une solution particulière $(x_0, y_0)$.
  \item Achever la résolution de l'équation $(E)$.
\end{enumerate}

% ====================== EXERCICE 2 ======================
\exo{2}{8 points}

Le plan est orienté. Soit $ABC$ un triangle équilatéral de sens direct, de centre de gravité $E$.

\begin{enumerate}
  \item Faire une figure. On prendra $AB = 4$ cm.
  \item Construire le cercle $(C)$ de centre $A$ passant par $B$.
  \item Construire le point $F$ symétrique de $A$ par rapport à $E$.
  \item Montrer que les droites $(CF)$ et $(CA)$ sont perpendiculaires.
\end{enumerate}

Soit $(\Gamma)$ l'hyperbole de cercle principal $(C)$ et dont une directrice est la droite $(BC)$.

\begin{enumerate}[start=5]
  \item Montrer que $F$ est un foyer de $(\Gamma)$. Préciser son axe focal.
  \item On désigne par $G$ le projeté orthogonal de $F$ sur la droite $(BC)$, et $H$ le point de l'axe focal tel que $\vv{AH} = 3\vv{AG}$.
  \begin{enumerate}[label=\alph*.]
    \item Construire $H$.
    \item On pose $AF = c$ et $AB = a$. Que représente le rapport $\dfrac{AF}{AB}$ pour $(\Gamma)$ ? Montrer que ce rapport est égal à $\dfrac{2\sqrt{3}}{3}$.
  \end{enumerate}
  \item Construire le point $I$ du plan tel que $\vv{GI} = \dfrac{c}{a}\vv{GH}$.
  \item Soit $(d)$ la perpendiculaire à l'axe focal passant par $H$. Construire le point $J$ de $(\Gamma)$ situé sur $(d)$ et situé dans le demi-plan délimité par la droite $(AH)$ contenant le point $C$.
  \item On note $S$ le sommet de $(\Gamma)$ associé au foyer $F$. Construire l'arc $(H)$ de $(\Gamma)$ d'extrémités $J$ et $S$.
  \item On désigne par $s$ la similitude plane indirecte définie par : $s = h \circ S_{(BC)}$ où $h$ est l'homothétie de centre $A$ et de rapport $\dfrac{1}{2}$ et $S_{(BC)}$ la symétrie orthogonale d'axe $(BC)$.
  \begin{enumerate}[label=\alph*.]
    \item Déterminer l'axe $(\Delta)$ et le centre $(\Omega)$ de $s$.
    \item Construire l'arc $(H')$, image de $(H)$ par $s$.
    \item Déterminer l'excentricité de $(\Gamma')$ image de $(\Gamma)$ par $s$.
  \end{enumerate}
\end{enumerate}

% ====================== EXERCICE 3 ======================
\exo{3}{5 points}

Soit $f$ la fonction numérique définie sur l'intervalle $I = [0\,;+\infty[$ par : $f(x) = (x-1)\ln(x+1)$. On pose pour tout $x \ge 0$, $F(x) = \displaystyle\int_1^x f(t)\,dt$.

\begin{enumerate}
  \item \begin{enumerate}[label=\alph*.]
    \item Prouver que $F$ est dérivable sur $I$ et que pour tout $x \in I$ ; $F'(x) = f(x)$.
    \item En déduire le sens de variation de $F$ sur $I$.
  \end{enumerate}
  \item On admet que pour tout $x \ge 2$, on a $f(x) \ge x - 1$.
  \begin{enumerate}[label=\alph*.]
    \item Prouver que pour tout $x \ge 2$, $F(x) \ge \dfrac{1}{2}(x-1)^2$.
    \note{Erreur dans l'énoncé : L'inégalité $F(x) \ge \dfrac{1}{2}(x-1)^2$ pour tout $x \ge 2$ est fausse. Substituer cette question par : Prouver que pour tout $x \ge 2$, $F(x) \ge F(2) + \dfrac{x^2}{2} - x$.}
    \item En déduire $\displaystyle\lim_{x\to+\infty} F(x)$ et $\displaystyle\lim_{x\to+\infty} \dfrac{F(x)}{x}$.
    \item Dresser le tableau de variation de $F$.
    \item Donner l'allure de la courbe $(C)$ représentant la fonction $F$ dans un repère orthonormal $(O,\vi,\vj)$ d'unité graphique 2 cm. On donne $F(0) = 0{,}13$ ; $F(2) = 0{,}5$.
  \end{enumerate}
  \item Soit la suite numérique $(u_n)$ définie pour tout entier naturel $n$ par : $u_n = \displaystyle\int_n^{n+1} f(t)\,dt$.
  \begin{enumerate}[label=\alph*.]
    \item Vérifier que $u_n = F(n+1) - F(n)$.
    \item En utilisant le théorème des inégalités des accroissements finis sur l'intervalle $[n\,;n+1]$, montrer que pour tout $n \in \N$, on a $f(n) \le u_n \le f(n+1)$.
    \note{Erreur dans l'énoncé : l'encadrement $f(n) \le u_n \le f(n+1)$ n'est pas vérifié pour tout $n \in \N$ mais pour tout $n \in \N^*$.}
  \end{enumerate}
\end{enumerate}

% ====================== EXERCICE 4 ======================
\exo{4}{3 points}

Jean s'amuse régulièrement sur un terrain de football avec le gardien de but. L'épreuve consiste à tirer au but et à observer le résultat obtenu. On admet que :
\begin{itemize}
  \item la probabilité que Jean réussisse le premier tir au but est de 0,7 ;
  \item s'il réussit le premier tir, alors la probabilité de réussir le second est de 0,8 ;
  \item s'il manque le premier tir, la probabilité de réussir le second est de 0,4.
\end{itemize}
On note $R_1$ l'événement « le premier tir au but est réussi » et $R_2$ l'événement « le second tir au but est réussi ».

\begin{enumerate}
  \item Construire l'arbre de probabilité correspondant à cette expérience.
  \item Calculer la probabilité pour que les deux tirs au but soient réussis.
  \item Calculer la probabilité pour que le second tir au but soit réussi.
  \item On note $A$, l'événement « Jean a réussi exactement un tir au but ». Calculer $P(A)$.
\end{enumerate}

\end{document}
